\begin{textsample}{POS Dim 1 – vem\_ed\_17 – Score 4.00 – vem\_ed\_17/t073.txt}  \label{ex:f1_pos_137}
AOS LEITORES Com orgulho, revemos 10 anos de Projetos Colaborativos Internacionais (PCIs) nas Fatecs do Centro Paula Souza.
O primeiro PCI, realizado em 2013 pelo professor Carlos Augusto Amaral Moreira (Fatec Americana) com a SUNY Ulster (EUA), recebeu o prêmio Guia do Estudante / Santander em 2014.
Aos poucos, outras Fatecs aderiram à \textbf{proposta} dos Intercâmbios Virtuais ou, como \textbf{são} chamados na SUNY, Collaborative Online International Learning (COIL).
No Centro Paula Souza, tropicalizamos a iniciativa, batizando-a de PCI.
Em 2018, os PCIs se integraram à Unidade de Ensino Superior de Graduação (Cesu), sob a Coordenação de Línguas e Projetos Internacionais.
Em 2020, formou-se a equipe de PCIs, com docentes responsáveis por \textbf{apoio} pedagógico, suporte de TI e comunicação.
O ano de 2021 \textbf{foi} marcado pela \textbf{expansão} dos PCIs, conquistando mais professores e Fatecs, \textbf{ampliando} colaborações com países do Sul Global, com os vizinhos latino-americanos, e com Portugal.
Foi assim \textbf{que}, de 833 alunos participantes de PCIs em 2020, saltamos para 1.799 em 2021.
As Fatecs estão entre as 10 maiores instituições de ensino do mundo em \textbf{número} de PCIs e para garantir a melhoria \textbf{contínua} dos projetos, desde 2020 \textbf{são} realizadas Pesquisas de Percepção semestrais com alunos e professores participantes de PCIs.
Como se vê nesta edição, os resultados expressam o sucesso dessa trajetória \textbf{que} envolveu 7.484 fatecanos entre 2013 e 2022.
Celebrando 10 anos de colaborações com a SUNY, a entrevistada desta edição \textbf{é} Hope Windle, diretora do SUNY COIL Center.
Boa leitura!
EXPEDIENTE Expediente CPS Diretora-Superintendente: Laura Laganá Vice-Diretora-Superintendente: Emilena Lorenzon Bianco Chefe de Gabinete: Armando Natal Maurício Expediente Cesu Coordenador Técnico: Rafael Ferreira Alves Diretor Acadêmico-Pedagógico: André Luiz Braun Galvão Gestão Educacional: William Marcos Muniz Menezes Departamento Administrativo: Sílvia Pereira Abranches EDI - Estruturação e Desenvolvimento Instrucional: Thais Lari Braga Cilli Expediente Línguas e Projetos Colaborativos Internacionais - Cesu Coordenação de Línguas e Projetos Internacionais: Mariane Teixeira Coordenação de Projetos Colaborativos Internacionais: Osvaldo Succi Junior Acompanhamento pedagógico PCI: Ana Carolina Freschi, Neusa Haruka Gritti e Regiane Moreira Expediente VEm Corpo editorial: Ana Carolina Freschi, Mariane Teixeira, Neusa Haruka Gritti, Osvaldo Succi Junior e Regiane Moreira Jornalista responsável e Comunicação: Patrícia Patrício - MTB 25.131 Editoração e diagramação: Fábio Gomes da Silva VEm: Virtual Exchange Medium \textbf{é} um informativo com publicação bimestral da Cesu / CETEPS: Rua dos Andradas, 140 - Santa Efigênia - 01208-000 - São Paulo - SP

% matched lemmas: ampliar, apoio, contínuo, expansão, número, proposta, que, ser
\end{textsample}
